\section{Introduction}\label{sec:intro}
The three-dimensional incompressible Navier--Stokes equations with
viscosity $\nu>0$ and external force $f$ are
\begin{equation}\label{eq:NS}
 \partial_tu+(u\cdot\nabla)u-\nu\Delta u+\nabla p=f,
 \qquad \nabla\cdot u=0,\qquad u(\cdot,0)=a.
\end{equation}
We work on $\R^3$ and on the unit torus $\TT=\R^3/\mathbb Z^3$.
For fixed $a$ and $T>0$, we ask in which force norms the smooth forces
producing classical breakdown by $T$ are dense.

The geometric argument is a divergence-free version of the smooth Urysohn
construction: make a given smooth velocity vanish near a small compact set,
leave it unchanged outside a slightly larger set, and place a compact blowup
solution in the resulting region. Smooth cutoffs are supplied by the classical
smooth Urysohn lemma~\cite[Corollary~1.3.4]{petersen}. To preserve
incompressibility, apply the cutoff to a local vector potential rather than
to the velocity itself. The remaining equation error is absorbed into the
external force. Our blowup building block is the following result of
OpenAI~\cite[Theorem~1.1]{openai}.

\Needspace{12\baselineskip}
\begin{theorem}\label{thm:packet}
For every $\nu>0$ there exist a force
$f\in C_c^\infty(\R^3\times(0,\infty);\R^3)$, a compact set
$K\subset\R^3$, and smooth velocity and pressure fields $u,p$ on
$\R^3\times[0,1)$ satisfying
\begin{equation}\label{eq:packet}
 \partial_tu+(u\cdot\nabla)u-\nu\Delta u+\nabla p=f,
 \qquad \nabla\cdot u=0,\qquad u(\cdot,0)=0,
\end{equation}
such that $\supp u(\cdot,t)\cup\supp p(\cdot,t)\subset K$ for every
$0\le t<1$,
\begin{equation}\label{eq:packetblowup}
 \sup_{0\le t<1}\norm{u(t)}_{L^2(\R^3)}<\infty,
 \qquad \limsup_{t\uparrow1}\norm{u(t)}_{L^\infty(\R^3)}=\infty.
\end{equation}
Consequently, there is no smooth solution $(u,P)$ on
$\R^3\times[0,\infty)$ with the same force and initial datum whose
kinetic energy is uniformly bounded
$\sup_{t\ge0}\frac12\int_{\R^3}|u(x,t)|^2\dd x<\infty$.
\end{theorem}

Fix one such building block $(U,P,F)$~\cite{openai,lean}.
The local gluing strategy parallels
Enciso, Pe\~nafiel-Tom\'as and Peralta-Salas~\cite[Theorem~1.5,
Lemma~2.11, Sections~11.1 and~11.3]{eppt2025}: first modify a given smooth
solution locally, then glue in a blowup building block. Their construction
produces weak unforced Euler solutions by removing Reynolds stress with
convex integration. Here the force is allowed to change, and the glued
solution is classical before its blowup time.

More explicitly, for a given smooth velocity $v=\nabla\times A$ locally,
choose smooth cutoffs and set
\[
 w_\eps=-\nabla\times(\eta_\eps\theta_\eps A),\qquad
 u_\eps=v+w_\eps+U_\eps.
\]
The corrected velocity $v+w_\eps$ vanishes near the support of the rescaled
building block $U_\eps$. Hence both mixed nonlinear terms vanish. The force
is the sum of the given force, the rescaled force, and a smooth cutoff
correction. Plancherel's theorem and the Fourier dilation formula, followed
by the time change $t=t_\eps+\eps^2\tau$, give the force factor
$\eps^{2/q-3/2-s}$ in $L^q_t\dot H^s_x$, wherever finite.
For $0\le s\le1$, the corresponding fractional estimate also follows from
the classical Sobolev interpolation inequality
\[
 \norm{z}_{\dot H^s}\le\norm{z}_2^{1-s}\norm{\nabla z}_2^s,
\]
which is H\"older's inequality on the Fourier side; see
\cite[Chapter~4, Sections~1--2]{taylor2011}. The cutoff correction has one
additional power of $\eps$. Negative orders require the separate
low-frequency estimate in Section~\ref{sec:whole}.

The resulting sharp thresholds for zero initial velocity are
\begin{center}
\begin{tabular}{lll}
\toprule Domain & Force topology & Density holds exactly when\\
\midrule
$\TT$ & $L^1_tH^s_x$ & $s<1/2$\\
$\R^3$ & $L^1_tH^s_x$ & $s<1/2$\\
$\R^3$ & $L^2_tH^s_x$ & $s<-1/2$\\
\bottomrule
\end{tabular}
\end{center}
Below threshold, density holds for every fixed admissible smooth initial
velocity. Near a solution smooth through $T$, the approximation preserves
its initial velocity and earlier history, blows up exactly at $T$, and
converges in energy and dissipation. For a general given force the conclusion
is breakdown by $T$, since a solution that already breaks down earlier needs
no modification. Non-density at and above threshold follows from the forced
critical regularity estimates, not from scaling alone.

For comparison, Hofmanov\'a, Zhu and Zhu~\cite[Theorem~1.4, Corollary~4.6,
Theorem~5.1]{hzz2024} obtain force-density results for nonunique Leray--Hopf
solutions. Our property is classical breakdown under a smooth force.
Classical local theory and the critical velocity framework are supplied by
Tao~\cite{tao2013,tao2018}. Section~\ref{sec:preliminaries} fixes conventions and collects
the standard estimates used below; Sections~\ref{sec:torus}
and~\ref{sec:whole} give the periodic and whole-space arguments.
